\part{Arithmetic of Integers}

\section{September 16th: Division}

\begin{defbox}
   Let $a$, $b \in \ZZ$. We say that \textbf{$a$ divides $b$}, denoted $a \mid b$, if there exists $q \in \ZZ$ such that $b = aq$.
\end{defbox}
\begin{propbox}
   Let $a$, $b$, $c \in \ZZ$. The following statements hold:
   \begin{enumerate}
      \item $a \mid b \implies a \mid -b$.
      \item $a \mid b \implies a \mid bc$.
      \item $a \mid b$ and $a \mid c \implies a \mid (b \pm c)$.
      \item $a \mid b$ and $b \mid c \implies a \mid c$.
   \end{enumerate}
\end{propbox}
The proofs of each of these statements are very easy by definition.

\begin{thmbox}[Quotient and Remainder Theorem]
   Let $a$, $b \in \ZZ$ where $b \neq 0$. Then there exists $q$, $r \in \ZZ$ such that $a = bq + r$ and $0 \leq r < |b|$. 
\end{thmbox}
\begin{proof}
   We must show existence and uniqueness of $q$ and $r$ in this proof. We begin by showing that $q$ and $r$ exist. Define:
   \begin{align*}
      S &:= \{r \in \NN \colon r =a - bq, q \in \ZZ\} \\
        &=\{a - bq \colon q \in \ZZ, a - bq \geq 0\}
   \end{align*}
   Then $S \subseteq \NN$ and $S \neq \emptyset$. Indeed, if $a \geq 0$, 
\end{proof}
